%==============================================================================
% IMPORTANT: Every numerical table and figure in this file must be regenerated
% after changing sensing_min_power_fraction. Do not use legacy zero-lower-bound
% results as evidence for the participation-constrained model.
\section{Numerical Results}\label{sec:numerical}
%==============================================================================

\subsection{Simulation Setup}
Unless otherwise stated, the Cell-Free ISAC network contains $M=6$ APs with
$N_t=2$ antennas per AP, $K=3$ UEs, and $P=2$ targets. Each target state has
dimension $N_\theta=2$, the deployment area is $400\,\mathrm{m}\times
400\,\mathrm{m}$, and the AP power limit is $20$ dBm ($0.1$ W). The nominal
CSI uncertainty radius is $\epsilon_h=0.05$, the communication and sensing
SINR thresholds are $0$ dB, and the automatically calibrated PCRB allowance
uses $\Gamma_{\mathrm{alpha}}=3$. A result is counted as feasible only after
binary-topology recovery, fixed-$b$ re-optimization, and physical validation.

\begin{table}[htbp]
\centering
\caption{Main simulation configuration.}
\label{tab:simulation-setup}
\begin{tabular}{p{0.56\linewidth}p{0.30\linewidth}}
\toprule
Parameter & Value\\
\midrule
APs and antennas per AP & $M=6$, $N_t=2$\\
UEs, targets, and state dimension & $K=3$, $P=2$, $N_\theta=2$\\
Required sensing APs per target & $N_{\mathrm{req}}=3$\\
Deployment area and AP power limit & $400\,\mathrm{m}\times400\,\mathrm{m}$, $0.1$ W\\
CSI uncertainty radius & $\epsilon_h=0.05$\\
Communication and sensing SINR targets & $0$ dB, $0$ dB\\
PCRB allowance & Automatic calibration, $\Gamma_{\mathrm{alpha}}=3$\\
Dedicated sensing participation floor & $P_{\min}^{\rm sen}=0.01P_{\max}=1$ mW\\
Solver and implementation & MATLAB R2024b, CVX, MOSEK 11.2.2\\
\bottomrule
\end{tabular}
\end{table}

\subsection{Binary-DC Recovery and Topology Quality}
The penalty ablation is evaluated over $25$ common scenarios. Relative to the
unpenalized continuous relaxation, the median binary residual decreases from
$4.274\times10^{-1}$ to $5.940\times10^{-5}$ when the binary DC penalty is
enabled. The rank residual remains below $10^{-8}$ for every ablation mode.
In this MU--MISO configuration, the transmit covariance is already nearly
rank one after SDR; consequently, the rank penalty is retained as a general
formulation component but is not claimed to yield an additional empirical
gain in this setting. Every reported physical point is obtained only after
binary topology recovery, fixed-$b$ re-optimization, and explicit validation.

\begin{table}[htbp]
\centering
\caption{Current-model topology and recovery evidence over $30$ common scenarios.}
\label{tab:main-mc}
\begin{tabular}{p{0.56\linewidth}p{0.30\linewidth}}
\toprule
Metric & Result\\
\midrule
Binary residual, unpenalized / dual-DC median & $4.268\times10^{-1}$ / $6.229\times10^{-5}$\\
FIM fixed-topology feasibility & $30/30$\\
FIM versus nearest-AP mean power reduction ($N_{\rm req}=3$) & $21.27\%$\\
FIM versus random: jointly feasible cases ($N_{\rm req}=3$) & $10/10$; mean reduction $73.30\%$\\
FIM power gap to SDR, median & $9.73\%$\\
\bottomrule
\end{tabular}
\end{table}

Fig.~\ref{fig:architecture} depicts the asymmetric signaling architecture.
The communication covariances $\{\mat{W}_k\}$ remain globally cooperative,
whereas $b_{mp}$ requires each selected AP to allocate nonzero dedicated
sensing covariance $\mat{S}_p$ to target $p$. The topology comparison shows
that the FIM-aware candidate is not interchangeable with a path-loss-only or
random association: at the main setting $N_{\rm req}=3$, its mean power is
$21.27\%$ below the nearest-AP baseline, while random association is feasible
in only $10$ of $30$ trials. Support
stability is used only as an engineering stopping rule; final binary
feasibility is certified by fixed-$b$ re-optimization.

\begin{figure*}[htbp]
\centering
\includegraphics[width=\textwidth]{../experiment_packages/v1.0/figures/fig2_dual_dc_ablation.png}
\caption{Penalty ablation over $25$ common scenarios. The binary DC penalty,
rather than the rank penalty, drives the association relaxation to the binary
boundary in the present MISO setting. Rank residuals are already negligible
after SDR, while all modes use the same fixed-$b$ physical validation step.}
\label{fig:dual-dc-ablation}
\end{figure*}

\begin{figure}[htbp]
\centering
\includegraphics[width=\linewidth]{../experiment_packages/v1.0/figures/fig1_system_architecture.png}
\caption{Cell-Free ISAC architecture. Global communication beamforming is not
gated by $b_{mp}$, whereas each selected AP--target pair radiates a
nonzero target-specific dedicated sensing component.}
\label{fig:architecture}
\end{figure}

\subsection{Cluster Size, QoS, and Recovery}
Fig.~\ref{fig:cluster-size} uses $30$ common seeds for every
$N_{\mathrm{req}}\in\{2,3,4,5,6\}$. The mean power decreases from
$37.65$ mW at $N_{\mathrm{req}}=2$ to $29.50$ mW at
$N_{\mathrm{req}}=4$, then rises to $30.94$ mW at $N_{\mathrm{req}}=6$.
Thus, the extra geometric diversity initially reduces tracking cost, whereas
the mandatory per-pair sensing floor eventually offsets this gain. All tested
points are physically feasible. The QoS
statistics in Fig.~\ref{fig:qos-cluster} show that the PCRB constraint is
numerically active: the maximum target PCRB ratio is essentially one. The
worst-user nominal communication and sensing-SINR margins remain
nonnegative. A recomputation from the returned covariance matrices gives a
maximum PCRB ratio numerically equal to one at the displayed precision.

\begin{figure}[htbp]
\centering
\includegraphics[width=\linewidth]{../experiment_packages/v1.0/figures/fig3_cluster_size_tradeoff.png}
\caption{Physical feasibility, total transmit power, and runtime versus
the required sensing-cluster size. Error bars show the 10th--90th
percentile interval.}
\label{fig:cluster-size}
\end{figure}

\begin{figure}[htbp]
\centering
\includegraphics[width=\linewidth]{../experiment_packages/v1.0/figures/fig4_qos_vs_cluster_size.png}
\caption{QoS statistics versus sensing-cluster size. Shaded regions denote
the 10th--90th percentile interval across feasible trials.}
\label{fig:qos-cluster}
\end{figure}

\subsection{Statistical Double-DC Stabilization}
To distinguish continuous-phase stabilization from final physical recovery,
we run the dual-DC SCA for six fixed iterations on 25 common seeds and include
the unpenalized SDR point as iteration zero.  The discrete candidate recovery
stage is intentionally omitted in this diagnostic, so the figure does not
replace the fixed-$b$ physical validation used for every reported solution.
The median maximum binary distance falls from $4.336\times 10^{-1}$ at the
SDR point to $1.139\times 10^{-3}$ after the first binary-DC step and to
$6.014\times10^{-5}$ after the second step.  All 25 Top-$N_{\rm req}$
supports are unchanged and recovery-ready from the first DC step onward,
while the median rank residual remains below $5\times10^{-9}$.  The
continuous power settles near $29.49$ mW by the second step.  Since the binary
penalty is continued, this power trace is a stabilization diagnostic rather
than a monotone descent claim.

\begin{figure}[htbp]
\centering
\includegraphics[width=\linewidth]{../experiment_packages/v1.0/figures/fig5_statistical_double_dc_convergence.png}
\caption{Statistical double-DC SCA stabilization over 25 common seeds. The
iteration-zero point is the unpenalized SDR relaxation. The recovery-ready
rate requires small rank residual, unchanged Top-$N_{\rm req}$ support, and
binary distance below the recovery threshold; it does not itself certify the
final physical solution.}
\label{fig:statistical-dc-stabilization}
\end{figure}

\subsection{Comparison with Alternative Association Rules}
Fig.~\ref{fig:method-comparison} compares the proposed binary-DC recovery
with an FIM-greedy topology, a nearest-AP topology, and a random cardinality-
feasible topology over $30$ common scenarios for each value of
$N_{\mathrm{req}}$. After a topology is selected, all physical baselines use
the same fixed-assignment robust beamforming and dedicated sensing-covariance
optimization. Hence, the comparison isolates the association and recovery
mechanism rather than attributing a continuous beamforming advantage to a
discrete rule. The SDR is retained only as a power lower bound and is not
included as a physical curve.

The proposed method is physically feasible in all $30$ trials at every
$N_{\mathrm{req}}$. The FIM-greedy rule is also feasible in all trials, but
its mean power is slightly higher. In contrast, the nearest-AP rule is
feasible in only $24/30$ trials at $N_{\mathrm{req}}=2$, while the random
rule is feasible in only $1/30$ trials. At $N_{\mathrm{req}}=3$, the proposed
method uses $30.83$ mW on average, compared with $30.93$ mW for FIM-greedy
and $39.29$ mW for nearest-AP association. The random baseline must be read
conditional on feasibility: it is feasible in only $10/30$ trials and has a
conditional mean power of $115.82$ mW. The average normalized PCRB remains
essentially one for every feasible method, confirming that the power gains do
not result from relaxed tracking quality. The reported sum rate is the
nominal aggregate spectral efficiency of the returned physical covariance;
it is descriptive rather than an additional optimization objective.

For every feasible physical solution, the number of nonzero AP--target
dedicated-sensing pairs equals $PN_{\mathrm{req}}$, as required by the
positive participation floor. The final panel of Fig.~\ref{fig:method-comparison}
therefore reports the number of \emph{unique} APs with nonzero sensing power,
which captures AP reuse across targets. At $N_{\mathrm{req}}=6$, every AP is
assigned to every target and all association rules reduce to the same topology
and the same solution.

\begin{figure*}[htbp]
\centering
\includegraphics[width=\textwidth]{../experiment_packages/v1.0/figures/fig10_method_comparison_vs_nreq.png}
\caption{Thirty-seed comparison of physical association and recovery methods
versus the required sensing-cluster size. Power, sum rate, sensing SINR, and
the number of active sensing APs are conditional on physical feasibility;
shaded bands denote the 10th--90th percentile interval.}
\label{fig:method-comparison}
\end{figure*}

For the same $30$ scenarios, the complete recovery pool achieves a mean power
of $30.83$ mW, compared with $30.93$ mW for the FIM fixed topology. The small
difference indicates that FIM geometry is already an effective topology
selector in the present architecture; the continuous robust beamforming and
dedicated sensing-covariance optimization then supply the remaining gain.
Fig.~\ref{fig:dimension} confirms the expected runtime increase with total
antenna dimension.

\begin{figure}[htbp]
\centering
\includegraphics[width=.78\linewidth]{../experiment_packages/v1.0/figures/fig7_dimension_sensitivity.png}
\caption{Runtime and physical feasibility versus total antenna dimension.}
\label{fig:dimension}
\end{figure}

\subsection{Physical-Setting Scalability and Geometry Stress Tests}
To assess robustness beyond the nominal configuration, we conduct a
factorized physical-setting study with $30$ common seeds per point.  One
factor is changed at a time from the reference setting
$(M,N_t,K,P,N_{\rm req})=(6,2,3,2,3)$: the AP count, the antennas per AP,
the UE load, the target load, the deployment side length, and the per-AP
power budget are swept over the ranges shown in Fig.~\ref{fig:physical-factors}.
Two controlled geometries are additionally considered: an edge UE--target
co-location and a pair of closely spaced targets.  This yields $22$ physical
configurations and $660$ scenarios.  The proposed recovery, FIM-greedy
association, and nearest-AP association are evaluated on every identical
realization, followed in each case by the same fixed-assignment robust
beamforming re-optimization and physical validation.

At $N_t=1$, automatic $\Gamma_{\rm track}$ calibration rejects all $30$
records before optimization because the isotropic-reference FIM is
unobservable.  This is a reference-observability screening event, rather than
a physical-infeasibility outcome, and is excluded from the following physical
method comparison.  For the remaining $630$ automatically observable
scenarios, the proposed and FIM-greedy methods are physically feasible in all
cases, while nearest-AP association is feasible in $600/630$ cases.  The
maximum recomputed mean PCRB ratio is $1.000272$; this small relative excess
is consistent with the $10^{-5}$ absolute solver-validation tolerance for
the scenario-scaled PCRB constraint.  The minimum reported communication and
sensing-SINR margins are $0.255$ dB and $9.49\times10^{-4}$ dB, respectively.

Fig.~\ref{fig:physical-factors} shows that FIM-greedy is a strong geometric
baseline: its power is close to that of the proposed recovery across the
factor sweeps.  This is consistent with the role of the FIM in selecting a
high-quality sensing topology, after which continuous robust beamforming
optimizes the transmit covariance.  In contrast, nearest-AP association
ignores angular information diversity.  Over the AP-count sweep, the proposed
method reduces paired conditional power by $12.4\%$--$18.2\%$ relative to
nearest-AP association.  The corresponding paired reduction is $21.3\%$ at
$K=4$ and $19.3\%$ at $P=3$.

The controlled geometries in Fig.~\ref{fig:pressure-geometries} expose the
same mechanism more sharply.  In the edge co-location scenario, nearest-AP
association is feasible in $29/30$ trials and requires $49.6\%$ more paired
power than the proposed recovery.  When targets are crowded, it is feasible
in only $7/30$ trials; conditional on those common feasible trials, its power
is $81.4\%$ higher.  The proposed and FIM-greedy methods remain feasible in
all $30$ trials in both stress geometries.  These results demonstrate that
distance alone is not a reliable proxy for the FIM geometry needed by the
PCRB constraint, whereas the proposed binary-DC recovery reliably realizes a
geometry-aware sensing cluster together with a validated robust covariance.

\begin{figure*}[htbp]
\centering
\includegraphics[width=\textwidth]{../experiment_packages/v1.0/results/extended_physical_mc/fig11_extended_physical_factors.png}
\caption{Extended $30$-seed physical-setting study. Solid curves show
conditional mean transmit power and dashed curves show physical feasibility.
All methods use the same fixed-topology robust beamforming re-optimization
after association selection. The $N_t=1$ point is omitted because automatic
reference-FIM calibration rejects it before optimization.}
\label{fig:physical-factors}
\end{figure*}

\begin{figure*}[htbp]
\centering
\includegraphics[width=\textwidth]{../experiment_packages/v1.0/results/extended_physical_mc/fig12_pressure_geometries.png}
\caption{Physical feasibility and conditional mean transmit power in the
controlled edge co-location and crowded-target geometries over $30$ common
seeds.}
\label{fig:pressure-geometries}
\end{figure*}

\subsection{Higher-Dimensional Scalability Validation}
We further consider a larger Cell-Free network with $M=12$ APs, $N_t=2$
antennas per AP, $K=6$ UEs, and $P=3$ targets.  At $N_{\rm req}=3$, the
experiment uses eight common realizations, four isolated MATLAB workers, and
a $60$~s MOSEK limit per convex subproblem.  The proposed recovery is
physically feasible in all $8/8$ cases, whereas the FIM-greedy and nearest-AP
topologies are feasible in $6/8$ cases and random association in $1/8$ case.
Conditional power must be interpreted on the common feasible set: on the six
realizations where FIM-greedy is feasible, the proposed method requires
$37.10$~mW versus $37.88$~mW; on the six nearest-AP common-feasible cases,
it requires $38.72$~mW versus $94.89$~mW.  The mean normalized PCRB of the
eight proposed solutions is $0.999987$, confirming that this gain is obtained
at the tracking boundary.  The mean end-to-end runtime is $496.3$~s per
scenario, which quantifies the computational cost of the high-dimensional
robust recovery rather than implying real-time operation.

\begin{figure}[htbp]
\centering
\includegraphics[width=\linewidth]{../experiment_packages/v1.0/results/large_scale_algorithm_validation/M12_K6_P3_Nreq3_seed01to08_workers4_t60/fig11_m12_scalability_validation.png}
\caption{High-dimensional validation for $M=12$, $K=6$, $P=3$, and
$N_{\rm req}=3$ over eight common realizations. Power is conditional on
physical feasibility; runtime is averaged over all realizations.}
\label{fig:m12-scalability}
\end{figure}

To verify that the cardinality effect is not specific to a single high-dimensional
operating point, we additionally sweep $N_{\rm req}\in\{2,3,4,5\}$ on five
common realizations of the same $M=12$, $K=6$, $P=3$ system.  Figure~\ref{fig:m12-nreq}
reports the physical feasibility rate, the conditional mean power, the mean
end-to-end runtime, and the normalized PCRB.  The proposed method is feasible
in $4/5$ cases at $N_{\rm req}=2$ and in all five cases for $N_{\rm req}\geq3$.
Its conditional mean power decreases from $45.07$~mW at $N_{\rm req}=2$ to
$35.72$~mW at $N_{\rm req}=4$, and is $35.76$~mW at $N_{\rm req}=5$.
FIM-greedy achieves the same full feasibility only from $N_{\rm req}=4$ onward,
while nearest-AP is feasible in only $2/5$ cases at $N_{\rm req}=2$ and requires
$104.40$~mW at $N_{\rm req}=3$.  The PCRB ratios remain at their respective
tracking boundaries.  These results confirm that the cardinality trade-off
persists at larger dimension, while also making explicit that the robust
recovery incurs a substantial computational cost.

\begin{figure*}[htbp]
\centering
\includegraphics[width=\textwidth]{../experiment_packages/v1.0/results/large_scale_algorithm_validation/M12_K6_P3_nreq2_4_5_seed01to05_workers4_t60/fig12_m12_nreq_scalability.png}
\caption{High-dimensional sensing-cluster cardinality sweep for $M=12$,
$N_t=2$, $K=6$, and $P=3$ over five common realizations.  Power is conditioned
on physical feasibility, whereas runtime is averaged over all realizations.}
\label{fig:m12-nreq}
\end{figure*}

\subsection{QoS-Requirement Sensitivity and CSI Robustness}
Fig.~\ref{fig:tradeoff-mc} averages a $5\times5$ QoS grid over five common
seeds. It reports the conditional mean transmit power required by the proposed
method as the robust communication-SINR target and PCRB allowance scale
$\alpha$ are varied; a smaller $\alpha$ imposes a stricter tracking
requirement. The required power increases when the communication target is
raised or $\alpha$ is reduced. All $25$ tested grid points are physically
feasible. Hence, the figure characterizes QoS-requirement sensitivity over the
tested range, rather than a communication--sensing Pareto boundary or the
boundary of the full feasible region.

Fig.~\ref{fig:robustness} compares the proposed robust design with a nominal
CSI baseline that removes the S-procedure and is designed with
$\epsilon_h=0$. Both designs are tested under identical independently sampled
complex channel perturbations inside the uncertainty ball. For
$\epsilon_h\in\{0.02,0.05,0.08\}$, the robust design has zero empirical
system outage in the $100$ perturbation samples per seed over $30$ seeds,
whereas the mean nominal-design outage is $89.0\%$, $89.6\%$, and $91.9\%$,
respectively. At the uncertainty radius $0.08$, the median robust power premium is only
$5.12\%$.
This sampled result complements, but does not replace, the S-procedure
certificate for the prescribed uncertainty set.

Finally, a $30$-seed participation-floor sensitivity study verifies the
physical interpretation of $b_{mp}$. For each positive floor in
$P_{\min}^{\rm sen}/P_{\max}\in\{0.5\%,1\%,2\%,5\%\}$, all $30$ returned
solutions satisfy the corresponding per-AP--target lower bound with zero
measured violation. The mean power is $30.83$ mW at the selected $1\%$ floor,
whereas a $5\%$ floor raises it to $37.11$ mW. This supports $1\%$ as a
non-degenerate participation level with modest energy cost.

\begin{figure}[htbp]
\centering
\includegraphics[width=\linewidth]{../experiment_packages/v1.0/figures/fig8_statistical_tradeoff.png}
\caption{QoS-requirement sensitivity of the proposed method over five common
seeds: conditional mean transmit power and unconditional physical-feasibility
rate versus the robust communication-SINR target and PCRB allowance scale.
Smaller $\alpha$ corresponds to a stricter tracking requirement.}
\label{fig:tradeoff-mc}
\end{figure}

\begin{figure}[htbp]
\centering
\includegraphics[width=\linewidth]{../experiment_packages/v1.0/figures/fig9_csi_robustness.png}
\caption{Robust and nominal CSI designs under ball-bounded channel
perturbations.}
\label{fig:robustness}
\end{figure}
